
\item @@M0@@ est le
morphisme structural @@M1@@.
\end{enumerate}

\section*{4. Théorème de Riemann--Roch : cas d'une immersion fermée régulière}

\textbf{4.1.} Soit @@M2@@ une immersion fermée régulière. Le théorème de Riemann--Roch pour @@M3@@ est essentiellement la relation VII 4.4, qui permet de calculer les classes de Chern de @@M4@@ à partir des @@M5@@. Plus précisément, si on introduit le @@M6@@-anneau @@M7@@ défini en V 5.5, alors l'homomorphisme de @@M8@@ dans @@M9@@ déduit de @@M10@@ est un @@M11@@-homomorphisme d'après VII 2.8 et VII 4.3. Le théorème de Riemann--Roch se ramène alors à une formule analogue concernant l'homomorphisme (additif) @@M12@@ défini par @@M13@@ c'est-à-dire à une relation de nature formelle, valable pour tout @@M14@@-anneau. Nous ne développerons pas ici cette démonstration, afin d'éviter de faire encore du formalisme sur les @@M15@@-anneaux (cf. @@M16@@ I 1.5), et nous reprendrons la méthode de [1] qui utilisait les schémas éclatés (méthode déjà employée pour la démonstration de VII 4.3, et qui est aussi celle de la fin de @@M17@@).

Comme il a été indiqué dans 3.9, on peut se limiter au cas d'une immersion régulière de codimension constante @@M18@@ vérifiant la condition de VII 4.1, le cas d'une immersion régulière générale se déduisant ensuite de 3.8. On reprendra dans ce paragraphe les notations de VII 3.1. On a donc le carré cartésien :
@@M19@@
